\section{Introduction}\label{sec:introduction}
Quantitative economic models increasingly place several margins of adjustment inside the same dynamic decision problem. A household chooses consumption and portfolio risk while investing in its own preferences. Recursive utility makes continuation utility part of the policy-evaluation equation. A sophisticated consumer anticipates the decisions of future selves, while a firm invests in anticipation of its rivals' policies. These problems require a computational method that evaluates the relevant continuation values and determines whether feasible changes in behavior improve the stated economic objective. The motivating preference-formation problem of \citet{qi2025} brings the internal and external margins together.

This paper develops Neural Bellman Operators (NBO) around policy evaluation, feasible improvement, and economic verification. A neural critic represents continuation utility, and a feasible actor represents the controls entering the state dynamics. The critic evaluates a fixed policy using its boundary conditions and utility domain. The actor then improves its Hamiltonian with the critic and its derivative channels held fixed. Independent verification measures the remaining evaluation error and the gain from feasible deviations. These steps give the trained networks an economic interpretation and specify the accuracy of the policy actually implemented.

Our first contribution is a common evaluation and improvement formulation for the economic applications. In smooth problems, a verification theorem translates boundary error, fixed-policy residual error, and feasible-action error into value and policy-loss bounds. For nonsmooth problems, a positive-weight Bellman formulation supplies the monotonicity needed for viscosity selection. A guarded evaluation map enforces a declared nodal tolerance and records the corrections required to do so. Complete-action bounds use signed derivative enclosures and verified face domination, including the effects of stopping switches. The representation can therefore be chosen for numerical flexibility while the economic acceptance criterion is expressed in the model's values and deviations.

The second contribution carries policy-specific verification through to a dense nonlinear capital diffusion on its original unbounded state space. A feasible analytical schedule supplies an absolute comparison with the economic optimum. Curvature controls the benefit from changing consumption, and a spectral bound controls the propagation of those changes through the coupled production process. A paired payoff identity then measures the value of a fitted correction to the schedule. Explicit transfer, tail, arithmetic, and sampling terms turn this comparison into a policy-specific continuous-time bound. The comparison class includes the original adapted controls. The candidate controller is specified by its update rule, memory, and information structure, so that the evaluated policy and the implemented policy are the same mathematical object. A product-space argument shows that independent private randomization preserves the economic value. An observation construction implements the internal controller from exact state histories, and a separate payoff-transfer theorem extends implementation to finite dates and noisy state measurements.

The third contribution makes the role of the Bellman evaluation block testable. In the capital economy, strong concavity bounds Hamiltonian loss by costate error and actor suboptimality. A conditional-projection identity separates variation in rollout costates from the error of a learned predictor. A performance-difference theorem integrates the resulting decision losses along the candidate policy's state distribution and separates costate regression error, rollout bias, and actor suboptimality in a welfare-improvement inequality. These results identify the quantities through which critic learning can improve an actor. The numerical study accordingly compares NBO with direct policy optimization, affine policies, and critic ablations. Direct method differences are formed on common simulated paths, with each fitted candidate fixed before final testing. Comparisons report economic endpoints, candidate-generation work, and verification work; checkpoint frontiers expose the relation between accuracy and computation. The recorded initialization streams describe the fitted policies in the experiment, and the stated confidence event concerns their independently evaluated payoffs.

The original applications give these computational objects distinct economic meanings. In the preference model, value gradients determine consumption and the costly adjustment of risk aversion, while a cross derivative links portfolio demand to preference risk. The policywise cost comparison establishes value monotonicity in the adjustment cost, and a specified compensating transfer gives a welfare interpretation to numerical error. Recursive utility is evaluated within a declared invariant utility domain. For temporal selves, the actor satisfies a spike-deviation condition against two endogenous continuation values. In the capital game, the independent criterion is a complete unilateral best response against fixed rival policies. The main text states this application structure, and the supplement gives the full models, derivations, and numerical records. Each conclusion is attached to the economic domain and approximation for which it is established.

The capital study also supplies two economically interpretable comparisons. A finite population of initial balance sheets varies aggregate capital and cross-sectional dispersion, while separate state stresses assess more severe starting conditions. A proportional management cost gives an exact resource interpretation to a verified policy improvement under log utility. An independent scalar nonlinear HJB calculation evaluates learned feedback rules against a full-action numerical reference. Its domain and grid refinements make the quality of this low-dimensional reference observable; the high-dimensional continuous-time bound follows from its own theorem. Together, these exercises distinguish the value of a feasible learned correction, the incremental benefit of its construction method, and its remaining distance to the optimum.

